\begin{table}[t]
\centering
\small
\caption{Validation-locked protocol for the primary analysis. The diagnostic upper-bound envelope is analyzed separately.}
\label{tab:locked-protocol}
\begin{threeparttable}
\begin{tabularx}{\textwidth}{p{0.24\textwidth}X}
\toprule
Component & Protocol \\
\midrule
Split design & Each model--dataset condition is split by correctness into 60\% train, 20\% validation, and 20\% held-out test for five seeds. The seeds are overlapping resplits, not independent datasets. \\
Feature fitting & Standardization and PCA are fitted on training data only; validation and test examples use the train-fitted transforms. \\
Probe families & Balanced logistic reliability selectors compare the archived lower-dimensional baseline, PCA-hidden vectors, and their concatenation. The baseline is output-only in eight conditions and includes three scalar layer maxima in sixteen conditions. The full grid uses $C\in\{0.1,1,10\}$ and PCA dimensions $\{16,32,64\}$. \\
Validation selection & The primary configuration is selected by validation AUROC within family. Sensitivity analyses select by validation AURC or Risk@80\% and equalize the candidate budget. \\
Protected test reporting & AUROC, failure AUPRC, Risk@80\%, AURC, E-AURC, and failure capture are evaluated after validation selection. Risk@80\% is a test-ranking metric, not a transferred numerical threshold. \\
Primary estimand & The five seed results are averaged within each of the 24 model--dataset conditions; equal-condition macro averages are reported. \\
Dependence-aware uncertainty & A two-way bootstrap resamples model and dataset clusters independently for 50,000 replicates. Leave-one-dataset-out and leave-one-model-out estimates assess cluster sensitivity. \\
Diagnostic envelope & Broader best-probe, layer, PCA, contrastive, and sampling analyses are supporting diagnostics and do not replace the validation-locked estimate. \\
\bottomrule
\end{tabularx}
\begin{tablenotes}[flushleft]\footnotesize\item The design prevents test-guided feature-family selection but does not constitute a preregistered equivalence study.\end{tablenotes}
\end{threeparttable}
\end{table}
